% ------------------------------------------------------------------------------
% Formatvorlage f�r Masterarbeiten an der Ohm-Hochschule N�rnberg
% ------------------------------------------------------------------------------
%   erstellt von Stefan Macke, 24.04.2009
%   http://blog.stefan-macke.de

% Dokumentenkopf ---------------------------------------------------------------
%   Diese Vorlage basiert auf "scrreprt" aus dem koma-script.
% ------------------------------------------------------------------------------
\documentclass[
    12pt, % Schriftgr��e
    DIV10,
    ngerman, % f�r Umlaute, Silbentrennung etc.
    a4paper, % Papierformat
    oneside, % einseitiges Dokument
    numbers=noenddot, % ohne Punkt am Ende von Kapitelnummern
    titlepage, % es wird eine Titelseite verwendet
    parskip=half, % Abstand zwischen Abs�tzen (halbe Zeile)
    headings=normal, % Gr��e der �berschriften verkleinern
    listof=totoc, % Verzeichnisse im Inhaltsverzeichnis auff�hren
    bibliography=totoc, % Literaturverzeichnis im Inhaltsverzeichnis auff�hren
    index=totoc, % Index im Inhaltsverzeichnis auff�hren
    captions=tableheading, % Beschriftung von Tabellen unterhalb ausgeben
    final % Status des Dokuments (final/draft)
]{scrreprt}

% Meta-Informationen -----------------------------------------------------------
%   Informationen �ber das Dokument, wie z.B. Titel, Autor, Matrikelnr. etc
%   werden in der Datei Meta.tex definiert und k�nnen danach global
%   verwendet werden.
% ------------------------------------------------------------------------------
% Meta-Informationen -----------------------------------------------------------
%   Definition von globalen Parametern, die im gesamten Dokument verwendet
%   werden k�nnen (z.B auf dem Deckblatt etc.).
%
%   ACHTUNG: Wenn die Texte Umlaute oder ein Esszet enthalten, muss der folgende
%            Befehl bereits an dieser Stelle aktiviert werden:
%            \usepackage[latin1]{inputenc}
% ------------------------------------------------------------------------------
\usepackage[latin1]{inputenc}


\newcommand{\titel}{Konzeption und Implementierung eines Produktkonfigurators als Plugin f�r das Shopsystem xt:Commerce Veyton}
\newcommand{\art}{Bachelorarbeit}
\newcommand{\autor}{Andr� Eife}
\newcommand{\studiengang}{Informatik}
\newcommand{\matrikelnr}{46881}
\newcommand{\erstgutachter}{Prof. Dr. Michael Frank}
\newcommand{\zweitgutachter}{Prof. Dr.-Ing. Dieter Vyhnal}
\newcommand{\jahr}{2011}
\newcommand{\ort}{Leipzig}
\newcommand{\schule}{Hochschule f�r Technik, Wirtschaft und Kultur Leipzig}
\newcommand{\fakultaet}{Fakult�t Informatik, Mathematik und Naturwissenschaften}
% ben�tigte Packages -----------------------------------------------------------
%   LaTeX-Packages, die ben�tigt werden, sind in die Datei Packages.tex
%   "ausgelagert", um diese Vorlage m�glichst �bersichtlich zu halten.
% ------------------------------------------------------------------------------
% Anpassung des Seitenlayouts --------------------------------------------------
%   siehe Seitenstil.tex
% ------------------------------------------------------------------------------
\usepackage[
    automark, % Kapitelangaben in Kopfzeile automatisch erstellen
    headsepline, % Trennlinie unter Kopfzeile
    ilines % Trennlinie linksb�ndig ausrichten
]{scrpage2}

% Anpassung an Landessprache ---------------------------------------------------
\usepackage[ngerman]{babel}

% Umlaute ----------------------------------------------------------------------
%   Umlaute/Sonderzeichen wie ���� direkt im Quelltext verwenden (CodePage).
%   Erlaubt automatische Trennung von Worten mit Umlauten.
% ------------------------------------------------------------------------------
\usepackage[latin1]{inputenc}
\usepackage[T1]{fontenc}
\usepackage{textcomp} % Euro-Zeichen etc.

% Schrift ----------------------------------------------------------------------
\usepackage{lmodern} % bessere Fonts
\usepackage{relsize} % Schriftgr��e relativ festlegen

% Grafiken ---------------------------------------------------------------------
% Einbinden von JPG-Grafiken erm�glichen
\usepackage[dvips,final]{graphicx}
% hier liegen die Bilder des Dokuments
\graphicspath{{Bilder/}}

% Befehle aus AMSTeX f�r mathematische Symbole z.B. \boldsymbol \mathbb --------
\usepackage{amsmath,amsfonts}

% f�r Index-Ausgabe mit \printindex --------------------------------------------
\usepackage{makeidx}

% Einfache Definition der Zeilenabst�nde und Seitenr�nder etc. -----------------
\usepackage{setspace}
\usepackage{geometry}

%option acronym um abk�rzungen leichter zu erstellen
\usepackage[acronym=true, nonumberlist, toc]{glossaries}

% zum Umflie�en von Bildern ----------------------------------------------------
\usepackage{floatflt}

\usepackage[bottom, hang]{footmisc}
\renewcommand{\footnotemargin}{0.3cm}

% zum Einbinden von Programmcode -----------------------------------------------
\usepackage{listings}
\usepackage{xcolor} 
\definecolor{hellgelb}{rgb}{1,1,0.9}
\definecolor{colKeys}{rgb}{0,0,1}
\definecolor{colIdentifier}{rgb}{0,0,0}
\definecolor{colComments}{rgb}{1,0,0}
\definecolor{colString}{rgb}{0,0.5,0}
%\lstset{
%    float=htbp,
%    numbers=none,
%    basicstyle=\ttfamily\color{black}\small\smaller,
%    identifierstyle=\color{colIdentifier},
%    keywordstyle=\color{colKeys},
%    stringstyle=\color{colString},
%    commentstyle=\color{colComments},
%    columns=flexible,
%    tabsize=2,
%    numberstyle=false,
%    extendedchars=true,
%    showspaces=false,
%    showstringspaces=false,
%    breaklines=true,
%    breakautoindent=true,
%    captionpos=b,
%    xleftmargin=17pt
%    }
\lstset{
    float=htbp,
    basicstyle=\ttfamily\smaller,
    columns=flexible,
    tabsize=2,
    extendedchars=true,
    showspaces=false,
    showstringspaces=false,
    numbers=left,
    numberstyle=\tiny,
    breaklines=true,
    breakautoindent=true,
    keywordstyle=\ttfamily\smaller,
    captionpos=b,
    xleftmargin=10pt,
    xrightmargin=10pt,
}

% wichtig f�r korrekte Zitierweise ---------------------------------------------
%\usepackage[square]{natbib}

% PDF-Optionen -----------------------------------------------------------------
\usepackage[
    bookmarks,
    bookmarksopen=true,
    colorlinks=true,
% diese Farbdefinitionen zeichnen Links im PDF farblich aus
%    linkcolor=red, % einfache interne Verkn�pfungen
%    anchorcolor=black,% Ankertext
%    citecolor=blue, % Verweise auf Literaturverzeichniseintr�ge im Text
%    filecolor=magenta, % Verkn�pfungen, die lokale Dateien �ffnen
%    menucolor=red, % Acrobat-Men�punkte
%    urlcolor=cyan, 
% diese Farbdefinitionen sollten f�r den Druck verwendet werden (alles schwarz)
    linkcolor=black, % einfache interne Verkn�pfungen
    anchorcolor=black, % Ankertext
    citecolor=black, % Verweise auf Literaturverzeichniseintr�ge im Text
    filecolor=black, % Verkn�pfungen, die lokale Dateien �ffnen
    menucolor=black, % Acrobat-Men�punkte
    urlcolor=black, 
    plainpages=false, % zur korrekten Erstellung der Bookmarks
    pdfpagelabels, % zur korrekten Erstellung der Bookmarks
    hypertexnames=false, % zur korrekten Erstellung der Bookmarks
    linktocpage % Seitenzahlen anstatt Text im Inhaltsverzeichnis verlinken
]{hyperref}
% Befehle, die Umlaute ausgeben, f�hren zu Fehlern, wenn sie hyperref als Optionen �bergeben werden
\hypersetup{
    pdftitle={\titel},
    pdfauthor={\autor},
    pdfcreator={\autor},
    pdfsubject={\titel},
    pdfkeywords={\titel},
}

% fortlaufendes Durchnummerieren der Fu�noten ----------------------------------
\usepackage{chngcntr}

% f�r lange Tabellen -----------------------------------------------------------
\usepackage{longtable}
\usepackage{array}
\usepackage{ragged2e}
\usepackage{lscape}

% Spaltendefinition rechtsb�ndig mit definierter Breite ------------------------
\newcolumntype{w}[1]{>{\raggedleft\hspace{0pt}}p{#1}}

% Formatierung von Listen �ndern -----------------------------------------------
\usepackage{paralist}

% bei der Definition eigener Befehle ben�tigt
\usepackage{ifthen}

% definiert u.a. die Befehle \todo und \listoftodos
\usepackage{todonotes}

% sorgt daf�r, dass Leerzeichen hinter parameterlosen Makros nicht als Makroendezeichen interpretiert werden
\usepackage{xspace}

% beseitigt Warnungen zu floats mit KOMA-Script
\usepackage{scrhack} 

% Captions linksb�ndig
%\usepackage[justification=raggedright,singlelinecheck=false]{caption}


% Erstellen des Glossars und Abk�rzungsverzeichnisses --------------------------
%Glossareintraege

\newglossaryentry{glos:Backend}{name=Backend, description={Administrationsoberfl�che, �ber welche der Administrator den Dienst verwalten kann}}
\newglossaryentry{glos:Frontend}{name=Frontend, description={Benutzersicht auf eine Seite oder ein Programm}}
\newglossaryentry{glos:Framework}{name=Framework, description={Ein Programmiergrundger�st mit vorgegebener Architektur}}
\newglossaryentry{glos:Plugin}{name=Plugin, description={Ein Modul zur funktionalen Erweiterung von Software}}

%Akronymeintraege

\newacronym{PHP}{PHP}{PHP: Hypertext Preprocessor}
\newacronym{ADOdb}{ADOdb}{Active Data Objects DataBase}
\newacronym{MySQL}{MySQL}{My Sequential Query Language}
\newacronym{HTML}{HTML}{Hypertext Markup Language}
\newacronym{XML}{XML}{Extensible Markup Language}
\newacronym{CSS}{CSS}{Cascading Style Sheets}
\newacronym{XAMPP}{XAMPP}{Apache MySQL PHP Perl}
\newacronym{Ajax}{Ajax}{Asynchronous JavaScript and XML}
\newacronym{W3C}{W3C}{World Wide Web Consortium}
\newacronym{HTTP}{HTTP}{Hypertext Transfer Protocol}
\newacronym{DOM}{DOM}{Document Object Model}
\newacronym{DBMS}{DBMS}{Datenbankmanagementsystem}

% Akronyme, welche im Flie�text nicht ausgeschrieben werden sollen
\glsunset{XAMPP}
\glsunset{PHP}
\glsunset{MySQL}
\glsunset{Ajax}
\glsunset{HTTP}
\glsunset{HTML}
\glsunset{XML}
\glsunset{ADOdb}
\glsunset{DOM}
\glsunset{DBMS}
\glsunset{CSS}

\makeglossaries

% Erstellung eines Index -------------------------------------------------------
\makeindex

% Kopf- und Fu�zeilen, Seitenr�nder etc. ---------------------------------------
% Zeilenabstand 1,5 Zeilen -----------------------------------------------------
\onehalfspacing

% Seitenr�nder -----------------------------------------------------------------
\setlength{\topskip}{\ht\strutbox} % behebt Warnung von geometry
\geometry{paper=a4paper,left=35mm,right=35mm,top=30mm, bottom=50mm}

% Kopf- und Fu�zeilen ----------------------------------------------------------
\pagestyle{scrheadings}
% Kopf- und Fu�zeile auch auf Kapitelanfangsseiten
\renewcommand*{\chapterpagestyle}{scrheadings} 
% Schriftform der Kopfzeile
\renewcommand{\headfont}{\normalfont}

% Kopfzeile
%\ihead{\large{\textsc{\titel}}\\ \textit{\headmark}}
\chead{}
\ohead{\headmark}
%\ohead{\includegraphics[scale=0.15]{\logo}}
\setlength{\headheight}{21mm} % H�he der Kopfzeile
% Kopfzeile �ber den Text hinaus verbreitern
%\setheadwidth[0pt]{textwithmarginpar} 
\setheadsepline[text]{0.4pt} % Trennlinie unter Kopfzeile

% Fu�zeile
\ifoot{\copyright\ \autor}
\cfoot{}
\ofoot{\pagemark}

% sonstige typographische Einstellungen ----------------------------------------

% erzeugt ein wenig mehr Platz hinter einem Punkt
\frenchspacing 

% Schusterjungen und Hurenkinder vermeiden
\clubpenalty = 10000
\widowpenalty = 10000 
\displaywidowpenalty = 10000

% Quellcode-Ausgabe formatieren
\lstset{numbers=left, numberstyle=\tiny, numbersep=5pt, breaklines=true}
\lstset{emph={square}, emphstyle=\color{red}, emph={[2]root,base}, emphstyle={[2]\color{blue}}}

% Fu�noten fortlaufend durchnummerieren
\counterwithout{footnote}{chapter}


% eigene Definitionen f�r Silbentrennung
% Trennvorschl�ge im Text werden mit \" angegeben
% untrennbare W�rter und Ausnahmen von der normalen Trennung k�nnen in dieser
% Datei mittels \hyphenation definiert werden


% eigene LaTeX-Befehle
% Eigene Befehle und typographische Auszeichnungen f�r diese

% einfaches Wechseln der Schrift, z.B.: \changefont{cmss}{sbc}{n}
\newcommand{\changefont}[3]{\fontfamily{#1} \fontseries{#2} \fontshape{#3} \selectfont}

% Abk�rzungen mit korrektem Leerraum 
\newcommand{\ua}{\mbox{u.\,a.\ }}
\newcommand{\zB}{\mbox{z.\,B.\ }}
\newcommand{\dahe}{\mbox{d.\,h.\ }}
\newcommand{\Vgl}{Vgl.\ }
\newcommand{\bzw}{bzw.\ }
\newcommand{\evtl}{evtl.\ }

\newcommand{\abbildung}[1]{Abbildung~\ref{fig:#1}}

\newcommand{\bs}{$\backslash$}

% erzeugt ein Listenelement mit fetter �berschrift 
\newcommand{\itemd}[2]{\item{\textbf{#1}}\\{#2}}

% einige Befehle zum Zitieren --------------------------------------------------
\newcommand{\Zitat}[2][\empty]{\ifthenelse{\equal{#1}{\empty}}{\cite{#2}}{\cite[#1]{#2}}}

% zum Ausgeben von Autoren
\newcommand{\AutorName}[1]{\textsc{#1}}
\newcommand{\Autor}[1]{\AutorName{\citeauthor{#1}}}

\newcommand{\code}[1]{\texttt{#1}}

% Das eigentliche Dokument -----------------------------------------------------
%   Der eigentliche Inhalt des Dokuments beginnt hier. Die einzelnen Seiten
%   und Kapitel werden in eigene Dateien ausgelagert und hier nur inkludiert.
% ------------------------------------------------------------------------------
\begin{document}

% Ebenentiefe im Inhaltsverzeichnis und bei Numerierung
\setcounter{secnumdepth}{2}
\setcounter{tocdepth}{2}

% Deckblatt und Abstract ohne Seitenzahl
\ofoot{}
\thispagestyle{plain}
\begin{titlepage}

\begin{center}

\Large{\textbf{\titel}}\\[1.5ex]
\huge{\textbf{\art}}\\[1.5ex]
\Large{im Studiengang \studiengang}\\[4ex]

\includegraphics[scale=0.35]{Bilder/HTWK-Logo.pdf}\\[4ex]
\large{\schule}\\[1ex]
\normalsize{\fakultaet}\\[3ex]

\normalsize
\begin{tabular}{w{5.4cm}p{6cm}}\\
vorgelegt von:  & \quad \autor\\[1.2ex]
Studiengang: & \quad \studiengang\\[1.2ex]
Erstgutachter:  & \quad \erstgutachter\\[1.2ex]
Zweitgutachter: & \quad \zweitgutachter\\[3ex] 
\end{tabular}

\copyright\ \jahr\\[9ex]

\end{center}
\end{titlepage}

% \listoftodos
\section*{Zusammenfassung}
\label{sec:Zusammenfassung}

\section*{Abstract}
\label{sec:Abstract}

\ofoot{\pagemark}

% Seitennummerierung -----------------------------------------------------------
%   Vor dem Hauptteil werden die Seiten in gro�en r�mischen Ziffern 
%   nummeriert.
% ------------------------------------------------------------------------------
\pagenumbering{Roman}

\addchap{Eidesstattliche Versicherung}
Ich, \autor, versichere hiermit, dass ich diese Bachelorarbeit mit dem Thema
\begin{quote}
\textit{\titel}
\end{quote}
selbstst�ndig ohne Hilfe Dritter und ohne Benutzung anderer als der angegebenen Quellen und Hilfsmittel verfasst habe. Alle den benutzten Quellen w�rtlich oder sinngem�� entnommenen Stellen sind als solche einzeln kenntlich gemacht.

Diese Arbeit ist bislang keiner anderen Pr�fungsbeh�rde vorgelegt und auch nicht ver�ffentlicht worden.

Ich bin mir bewusst, dass eine falsche Erkl�rung rechtliche Folgen haben wird.\\[6ex]



\ort, \today 	\hfill  	\rule{5cm}{0.5pt}\\
\hbox{}				\hfill		\textsc{\autor} \hspace{1.2cm} \\







 % Selbst�ndigkeitserkl�rung

\tableofcontents % Inhaltsverzeichnis


% arabische Seitenzahlen im Hauptteil ------------------------------------------
\clearpage
\pagenumbering{arabic}

% die Inhaltskapitel werden in "Inhalt.tex" includiert -------------------------
% Einzelne Inhaltskapitel
\sloppy
\chapter{Einleitung}
\label{cha:Einleitung}

\section{Motivation}


\section{Ziel der Arbeit}


\section{Aufbau der Bachelorarbeit}

\chapter{Grundlagen}
\label{cha:Grundlagen}


\chapter{Analyse aktueller Technik und Anforderungen}
\label{cha:Analyse}
\chapter{Entwurf}
\label{cha:Entwurf}

\chapter{Implementierung}
\label{cha:Implementierung}

\chapter{Fazit}
\label{cha:Fazit}



% Glossar und Abk�rzungsverzeichnis
% \printglossary[title=Glossar]
\deftranslation[to=German]{Acronyms}{Abk�rzungsverzeichnis}
% \printglossary[type=\acronymtype,style=listdotted,title=Abk�rzungsverzeichnis] 

\listoffigures % Abbildungsverzeichnis
% \listoftables % Tabellenverzeichnis
\renewcommand{\lstlistlistingname}{Verzeichnis der Listings}
\lstlistoflistings % Listings-Verzeichnis


% Literaturverzeichnis ---------------------------------------------------------
%   Das Literaturverzeichnis wird aus der BibTeX-Datenbank "Bibliographie.bib"
%   erstellt.
% ------------------------------------------------------------------------------
\bibliography{Bibliographie} % Aufruf: BibTeX-Datei
\bibliographystyle{alphadin} % DIN-Stil des Literaturverzeichnisses


% Anhang -----------------------------------------------------------------------
%   Die Inhalte des Anhangs werden analog zu den Kapiteln inkludiert.
%   Dies geschieht in der Datei "Anhang.tex".
% ------------------------------------------------------------------------------
% \begin{appendix}
%     \clearpage
%     \pagenumbering{roman}
%     \chapter{Anhang}
%     \label{sec:Anhang}
%     % Rand der Aufz�hlungen in Tabellen anpassen
%     \setdefaultleftmargin{1em}{}{}{}{}{}
%     % einzelne Anh�nge

\section{Verwendete Software}
\label{sec:Verwendete Software}
Im Folgenden werden die Programme vorgestellt, welche zur Erstellung dieser Arbeit und Entwicklung des Produktkonfiguratorplugins benutzt wurden.

\begin{itemize}
	\itemd{Dia}{S�mtliche Diagramme wurden mit dem Zeichenprogram Dia 0.97.1 erstellt.}
	\itemd{PHPEdit}{Der gesamte Quellcode des Produktkonfiguratorplugins wurde mit dem PHP Editor PHPEdit in der Version 3.2.0.8595 geschrieben.}
	\itemd{XAMPP}{Zum Testen des Plugins auf einem lokalen Computer wurde \gls{XAMPP} verwendet. Dabei musste auf die Version 1.6.7. zur�ckgegriffen werden um Kompatibilit�tsprobleme mit xt:Commerce Veyton und dem ben�tigten IonCube Loader zu vermeiden.}
	\itemd{FileZilla}{F�r den Transfer der Daten auf den Webserver des Onlineshops wurde die Version 3.2.7.1 des FTP-Client FileZilla benutzt.}
	\itemd{Firefox}{Die Verwaltung des Shopsystems sowohl im \gls{glos:Frontend} als auch im \gls{glos:Backend} wurde in dem Browser Firefox erledigt. Das Backend von xt:Commerce Veyton ist dar�ber hinaus auch lediglich in Firefox fehlerfrei darstellbar.}
	\itemd{\LaTeX}{Zur Erstellung dieser Arbeit wurde {\LaTeX} verwendet. Als Distribution wurde MiKTeX 2.9 verwendet und als Editor die Software TeXnicCenter sowie TextMate.}
	\itemd{JabRef}{Zur Verwaltung des Literaturverzeichnisses dieser Arbeit wurde die Software JabRef 2.6 verwendet.}
\end{itemize}
%\section{Anleitung}
\label{sec:Anleitung}

Anleitung zur Installation und Benutzung des Plugins im Shopsystem xt:Commerce Veyton.

\subsection{Installation}
F�r die Installation des Produktkonfigurators kopieren Sie den mitgelieferten Ordner \code{pc\_product\_configurator} in den Ordner \code{plugins} im Hauptverzeichnis Ihres Onlineshops befindet. In der Datei \code{product.html}, welche sich im Pfad \code{templates/xt\_default/xtCore/pages/product/} Ihres Onlineshops befindet, muss zudem noch der Code \code{\{hook key=pc\_product\_configurator\}} vor der Zeile \code{<div class="'addtobasket"'>} eingef�gt werden.\\
Anschlie�end melden Sie sich im Backend Ihres Shops an und gehen zu dem Men�punkt Inhalte / Plugins / deinstallierte Plugins. Dort sehen Sie unter der Modulklasse "`products"' den Produktkonfigurator. Klicken Sie auf das gr�ne Symbol rechts und best�tigen Sie die Installation. Bei Erfolg der Installation erhalten Sie eine entsprechende Meldung, welche wiederum best�tigt werden muss.\\
Das Plugin ist nach der Installation zun�chst noch inaktiv. Um es zu aktivieren gehen Sie unter Inhalte / Plugins / installierte Plugins. Dort sollten Sie den Produktkonfigurator in der Liste sehen. Ist dies nicht der Fall klicken sie auf "`neu laden"'. Bearbeiten Sie nun den Produktkonfigurator indem Sie zum Beispiel doppelt auf die entsprechende Zeile klicken. Im neu erscheinenden Fenster setzen Sie einen Haken bei "`Status"' und speichern die Bearbeitung. Der Produktkonfigurator ist nun aktiv und kann von Ihnen genutzt werden.

\subsection{Generelle Funktionsweise}
Um ein grobes Verst�ndnis im Umgang mit dem Produktkonfigurator zu erlangen soll an dieser Stelle ein grober �berblick �ber die Funktions- und Ablaufweise gegeben werden.\\
Der Produktkonfigurator erm�glicht es Produkten Gruppen von Komponenten zuzuweisen. Jede zugewiesen Gruppe entspricht einem Auswahlfeld auf der jeweiligen Produktseite im Shop. Diese Auswahlfelder beinhalteten zun�chst die Standardkomponenten, sofern diese Gruppen eine definiert haben. Ist keine Standardkomponente f�r eine Gruppe definiert bedeutet das, dass die Gruppe optional ist. Zur Auswahl stehen jeweils alle Komponenten, welche die Gruppe beinhaltet. Um die Gruppen besser gliedern zu k�nnen werden sie in �bergeordnete Gruppen eingeteilt.\\ 
Der Ablauf besteht also darin zun�chst Hauptgruppen- bzw. Kategorien anzulegen. In diesen Hauptgruppen k�nnen Untergruppen angelegt werden. Dies entspricht den Gruppen, welche sp�ter Produkten direkt zugewiesen werden k�nnen. Untergruppen beinhalten dann die einzelnen Komponenten, von denen optional eine als Standardkomponente deklariert werden kann.\\
Im Folgenden werden die einzelnen Schritte und deren Durchf�hrung n�her erl�utert.

\subsection{Komponentengruppen verwalten}
F�r die Verwaltung Ihrer Komponentengruppen ist ein gleichnamiger Men�punkt in der Shopsektion vorhanden. Nach einem Klick auf diesen Men�punkt befinden Sie sich zun�chst in der Ansicht f�r die Hauptkomponentengruppen. Hier k�nnen Sie neue Hauptgruppen anlegen oder vorhandene Hauptgruppen bearbeiten bzw. l�schen. Ebenso k�nnen sie Hauptgruppen kopieren, wobei s�mtliche enthaltene Inhalte wie Untergruppen und Komponente mitkopiert werden.\\
Um zur Ansicht f�r die Untergruppen einer bestimmten Hauptgruppe zu gelangen klicken Sie entweder in der Bearbeitenansicht unten auf den Button "`Komponentenuntergruppen"' oder in der �bersicht rechts von der entsprechenden Gruppe auf das Ordnersymbol. Sie befinden sich daraufhin in einer �bersicht, welche alle Untergruppen der jeweiligen Hauptgruppe anzeigt. Auch hier k�nnen Sie neue Untergruppen anlegen oder vorhandene Untergruppen bearbeiten bzw. l�schen.  Eine Besonderheit stellt hierbei das Kopieren dar. Das einfache Kopieren kopiert die Untergruppe innerhalb der jeweiligen Hauptgruppe. Es dupliziert die Gruppe quasi an der aktuellen Position. Das Kopieren von Untergruppen in Hauptgruppen bewirkt hingegen, dass sich ein extra Fenster �ffnet. Dieses zeigt Ihnen eine Baumstruktur aller von Ihnen angelegten Haupt- und Untergruppen. Sie k�nnen hier die Hauptgruppen ausw�hlen, in welche die Untergruppe kopiert werden soll. Dabei werden wie beim normalen Kopieren s�mtliche enthaltene Komponente mitkopiert.

\subsection{Komponenten verwalten}
Um die einzelnen Komponenten einer Untergruppe zu bearbeiten lassen Sie sich zun�chst, wie beschrieben, die Untergruppen einer gew�nschten Gruppe anzeigen. In dieser Ansicht klicken Sie auf den Button "`Komponenten"'. Sie finden diesen Button erneut entweder unten in der Bearbeitenansicht   oder in der �bersicht rechts von der entsprechenden Gruppe als Bausteinsymbol. Nach einem Klick auf einen der beiden Buttons wird Ihnen die Komponenten�bersicht angezeigt. Hier sehen Sie alle Komponenten der jeweiligen Untergruppe. Sie k�nnen neue Komponenten anlegen oder die bereits bestehenden verwalten. Neben dem Bearbeiten und L�schen von Komponenten stehen Ihnen wie bei den Untergruppen 2 Arten von Kopieren zur Verf�gung. Das einfache Kopieren kopiert die Komponente innerhalb der eigenen Untergruppe, dupliziert diese also. Das Kopieren einer Komponente in Untergruppen bewirkt, dass sich ein Fenster mit einer Baumhierarchie Ihrer angelegten Gruppen und Untergruppen �ffnet. Hier k�nnen Sie alle Untergruppen anw�hlen, in welche die Komponente kopiert werden soll.\\
Ein weiterer wichtiger Punkt bei der Verwaltung der Komponenten stellt die Auswahl der Standardkomponenten dar. Sie repr�sentiert die Komponente, welche im Shop von Anfang an ausgew�hlt ist, also zur Standardausstattung des jeweiligen Produktes geh�rt. Gibt es in einer Gruppe keine Komponente, welche als Standardkomponente definiert ist, gilt die Gruppe als optional. Im Shop ist zun�chst bei dieser Gruppe keine Komponente gew�hlt, aber es kann optional eine der enthaltenen Komponenten ausgew�hlt werden.

\subsection{Einem Produkt Komponentengruppen zuweisen}
Jedem im Shop verf�gbaren Artikel k�nnen Komponentengruppen zugewiesen werden. Dabei sind immer konkrete Untergruppen, mit deren enthaltenen Komponenten, gemeint. Um dies zu bewerkstelligen lassen Sie sich im Backend Ihres Shops das jeweilige Produkt entweder in der Bearbeitenansicht oder in der �bersicht anzeigen. In beiden Ansichten befindet sich bei aktivierten Produktkonfigurator ein Button "`Komponentengruppen zuweisen"'. Dies �ffnet ein Fenster, in dem alle Komponentengruppen und deren Untergruppen hierarchisch dargestellt sind. Um eine Untergruppe dem Produkt zuzuweisen, setzen Sie an die jeweilige Gruppe einen Haken und best�tigen Sie Ihre Eingabe mit einem Klick auf "`speichern"'. Dem Produkt sind nun die jeweilig ausgew�hlten Gruppen zugewiesen und stehen dem Nutzer nun in Ihrem Shop zur Verf�gung. Sie k�nnen jederzeit die Zuweisung der Gruppen durch erneutes �ffnen des Fensters ver�ndern und somit Gruppen hinzuf�gen oder entfernen.\\
Damit in Ihrem Shop der Produktkonfigurator zu dem jeweiligen Produkt auch angezeigt wird, m�ssen Sie in den Eigenschaften des Produkts noch einen Haken bei "`Produktkonfigurator"' setzen. Damit ist der Konfigurator f�r das jeweilige Produkte aktiviert und kann in Ihrem Shop genutzt werden.
%\section{Dokumentation}
\label{sec:Dokumentation}

Dokumentation des Plugins mittels PHPDocumentor

% \end{appendix}

% Index ------------------------------------------------------------------------
%   Zum Erstellen eines Index, die folgende Zeile auskommentieren.
% ------------------------------------------------------------------------------
%\printindex

\end{document}
